% ------------------------------------------------------------------------
%
\begin{frame}{Termination}

  When defining our functional rewrite systems, we allowed programs to
  not terminate (see pages~\pageref{termination1}
  and~\pageref{termination2}).

  \bigskip

  We could actually have imposed some syntactic restrictions on
  recursive definitions in order to guarantee the termination of all
  functions. A well\hyp{}known class of such terminating functions
  makes exclusively use of a bridled form of recursion called
  \textbf{primitive recursion}.

  \bigskip

  Unfortunately, many useful functions do not fit, easily or at all,
  in this framework and, as a consequence, most functional languages
  leave to the programmers the responsibility to check the termination
  of their programs. For theoretical reasons, it is not possible to
  provide a general criterion for termination, but some rules exist
  that cover many usages.

\end{frame}

% ------------------------------------------------------------------------
%
\begin{frame}{Termination is undecidable}

  Here is an informal argument that shows that there is no predicate
  on functions that decides where their subject terminates on all
  inputs. By contradiction, let us assume that there exists~\(f\) such that
  \begin{equation*}
    f(g) =
    \begin{cases}
      \mathsf{true}  & \text{if $g$ terminates on all inputs},\\
      \mathsf{false} & \text{otherwise}.
    \end{cases}
  \end{equation*}
  Now let us define a function \(h\) as follows:
  \begin{equation*}
    h(g) =
    \begin{cases}
      \text{loops forever}            & \text{if $f(g)$},\\
      \text{terminates on all inputs} & \text{otherwise.}
    \end{cases}
  \end{equation*}
  In other words, \(h\) looks at \(f(g)\) and does the opposite of
  what \(f\) predicts about \(g\).

\end{frame}

% ------------------------------------------------------------------------
%
\begin{frame}{Termination is undecidable}

  Let us consider now \(h(h)\):
  \begin{itemize}

    \item If \(f(h)\) is true, then, by definition of~\(f\),
      \(h\)~terminates on all inputs, in particular
      \(h(h)\)~terminates. But, by definition of~\(h\), since \(f(h)\)
      holds, \(h(h)\) loops on all inputs, yielding a contradiction.

    \item If \(f(h)\) is false, then, by definition of~\(f\),
      \(h\)~does not terminate on all its inputs. But, by definition
      of~\(h\), since \(f(h)\) does not hold, \(h(h)\) terminates on
      all inputs, yielding another contradiction.

  \end{itemize}

\end{frame}

% ------------------------------------------------------------------------
%
\begin{frame}{Termination / Factorial}

  Here is the definition of the factorial function as a functional
  rewrite system:
  \begin{align*}
    \fun{fact} \, 0 & \rightarrow 1;\\
    \fun{fact} \, n & \rightarrow n \times \fun{fact}(n-1).
  \end{align*}

  \medskip

  Given a call \(\fun{fact}(n)\), with \(n \geqslant 0\), the rewrites
  yield a series of calls with decreasing arguments:
  \(\fun{fact}(n-1)\), \(\fun{fact}(n-2)\), etc. until
  \(\fun{fact}(0)\).

  \bigskip

  Termination here ensues from the \textbf{strictly decreasing
    arguments}:
  \begin{equation*}
    n > n-1 > n-2 > \dots > 0.
  \end{equation*}
  No infinite descent.

\end{frame}

% ------------------------------------------------------------------------
%
\begin{frame}{Termination / Ackermann's function}

Consider the following example where \(m,n \in \mathbb{N}\):
\begin{equation*}
\begin{array}{@{}r@{\;}l@{\;}l@{}}
\fun{ack}(0,n)     & \smashedrightarrow{\alpha} & n+1;\\
\fun{ack}(m+1,0)   & \smashedrightarrow{\beta} & \fun{ack}(m,1);\\
\fun{ack}(m+1,n+1) & \smashedrightarrow{\gamma}
                   & \fun{ack}(m,\fun{ack}(m+1,n)).
\end{array}
\end{equation*}
This is a simplified form of Ackermann's function, an early example of
a total computable function which is not primitive recursive. It makes
use of double recursion and two parameters to grow values as towers of
exponents, for example,
\begin{equation*}
  \fun{ack}(4,3) \twoheadrightarrow 2^{2^{65536}} - 3.
\end{equation*}
It is not obviously terminating, because if the first argument does
decrease, the second largely increases.

\end{frame}

% ------------------------------------------------------------------------
%
\begin{frame}{Termination / Ackermann's function}

Let us define a well\hyp{}founded ordering on pairs, called
\textbf{lexicographic order}. Let \((\succ_A)\) and \((\succ_B)\) be
well\hyp{}founded orders on the sets \(A\) and~\(B\). Then \({(\succ_{A
  \times B})}\) defined as follows on \(A \times B\) is
well\hyp{}founded:
\begin{equation*}
(a_0,b_0) \succ_{A \times B} (a_1,b_1) :\Leftrightarrow \text{\(a_0
    \succ_A a_1\) or (\(a_0 = a_1\) and \(b_0 \succ_B b_1\)).}
\end{equation*}
If \(A=B=\mathbb{N}\) then \((\succ_A) = (\succ_B) = (>)\). \\ For
example, \((\underline{1},7) \prec (\underline{2},3)\) and
\((1,\underline{7}) \succ (1,\underline{3})\).

\bigskip

To prove that \(\fun{ack}(m,n)\) terminates for all \(m,n \in
\mathbb{N}\), first, we must find a well\hyp{}founded order on the
calls \(\fun{ack}(m,n)\), that is, the calls must be totally ordered
without any infinite descending chain.

\bigskip

Here, a lexicographic order on \((m,n) \in \mathbb{N}^2\) extended to
\(\fun{ack}(m,n)\) works:
\begin{equation*}
\fun{ack}(a_0,b_0) \succ \fun{ack}(a_1,b_1) :\Leftrightarrow
\text{\(a_0 > a_1\) or (\(a_0 = a_1\) and \(b_0 > b_1\)).}
\end{equation*}

\end{frame}

% ------------------------------------------------------------------------
%
\begin{frame}{Termination / Ackermann's function}

  Clearly, \(\fun{ack}(0,0)\) is the minimum element.

  \bigskip

  Second, we must prove that \fun{ack} rewrites to smaller calls. We
  are only concerned with rules \(\beta\)~and~\(\gamma\). With the
  former, we have the inequality
  \begin{equation*}
    \fun{ack}(m+1,0) \succ \fun{ack}(m,1).
  \end{equation*}
  With rule~\(\gamma\), we have the following inequalities:
  \begin{equation*}
    \begin{array}{@{}r@{\;}l@{\;}l@{}}
      \fun{ack}(m+1,n+1) & \succ & \fun{ack}(m+1,n),\\
      \fun{ack}(m+1,n+1) & \succ & \fun{ack}(m,p),
    \end{array}
  \end{equation*}
  for all values~\(p\), in particular when \(\fun{ack}(m+1,n)
  \twoheadrightarrow p\).\hfill\(\Box\)

  \bigskip

  Implicitly, we used \[x \succ y \Rightarrow f(x) \succ f(y),\]
  that is, \textbf{evaluation is (strictly) order\hyp{}preserving}.

\end{frame}

% ------------------------------------------------------------------------
%
\begin{frame}{Termination / Dependency pairs and subterm order}

  In general, the approach often taken to prove termination consists
  in finding a well\hyp{}founded order \((\succ)\) on the rewritten
  terms, which is entailed by the rewrite relation \((\rightarrow)\),
  that is,
  \begin{equation*}
    x \rightarrow y  \Rightarrow  x \succ y.
  \end{equation*}

  \medskip

  As we have seen, we do not always can order the whole terms
  \(x\)~and~\(y\): only the calls in them to a given function are
  enough. These are called the \textbf{dependency pairs}. Calls to
  other functions can be assumed to be terminating.

  \bigskip

  With the factorial and Ackermann's function, we have seen how to
  deal with one or more integer arguments to these calls.

  \bigskip

  What about other data types, like stacks (sometimes called `lists')
  and trees?

\end{frame}

% ------------------------------------------------------------------------
%
\begin{frame}{Termination / Dependency pairs and subterm order}

  One well\hyp{}founded order for lists (stacks) is the
  \textbf{immediate subterm order}, satisfying
  \begin{equation*}
    \cons{x}{s} \succ s \quad \text{and} \quad \cons{x}{s} \succ x.
  \end{equation*}

  The immediate subterm order simply means that a list is `greater'
  than its first element and the rest of the list (called the
  \textbf{tail}). It is a special case of the \textbf{proper subterm
    order}, in which a structured value is `greater' than any of its
  components, wherever they are located.

  \bigskip

  These orders are convenient to prove the termination of calls to
  functions whose definitions are recursive in terms of substructures
  of their arguments: this is quite a common case.

  \bigskip

  For example, let us recall the function catenating a list to
  another:
  \begin{align*}
    \fun{cat}(        \el,t) &\xrightarrow{\alpha} t\\
    \fun{cat}(\cons{x}{s},t) &\xrightarrow{\smash{\beta}}
    \cons{x}{\fun{cat}(s,t)}
  \end{align*}
  We simply have: \(\cons{x}{s} \succ s \Rightarrow
  \fun{cat}(\cons{x}{s},t) \succ \fun{cat}(s,t).\)

\end{frame}

% ------------------------------------------------------------------------
%
\begin{frame}{Termination / Insertion sort}

  Let us consider \textbf{insertion sort}:
\begin{equation*}
\begin{array}{@{}r@{\;}l@{\;}lr@{\;}l@{\;}l@{}}
  \fun{isort}\,\el
& \smashedrightarrow{\alpha}
& \el;
& \fun{ins}(x, \cons{y}{s})
& \smashedrightarrow{\gamma}
& \cons{y}{\fun{ins}(x,s)}, \,\text{if \(x \succ y\)};\\
  \fun{isort}(\cons{x}{s})
& \smashedrightarrow{\beta}
& \fun{ins}(x,\fun{isort}\,s).
&  \fun{ins}(x,s)
& \smashedrightarrow{\delta}
& \cons{x}{s}.
\end{array}
\end{equation*}
Let us consider a short example:
\begin{equation*}
\begin{array}{r@{\;}l@{\;}l}
\fun{isort}\,[3,1,2]
& \smashedrightarrow{\beta} & \fun{ins}(3,\fun{isort}\,[1;2])\\
& \smashedrightarrow{\beta}
& \fun{ins}(3,\fun{ins}(1,\fun{isort}\,[2]))\\
& \smashedrightarrow{\beta}
& \fun{ins}(3,\fun{ins}(1,\fun{ins}(2,\fun{isort}\,\el)))\\
& \smashedrightarrow{\alpha}
& \fun{ins}(3,\fun{ins}(1,\fun{ins}(2,\el)))\\
& \smashedrightarrow{\delta}
& \fun{ins}(3,\fun{ins}(1,[2]))\\
& \smashedrightarrow{\delta}
& \fun{ins}(3,[1;2])\\
& \smashedrightarrow{\gamma}
& \cons{1}{\fun{ins}(3,[2])}\\
& \smashedrightarrow{\gamma}
& \cons{1}{\cons{2}{\fun{ins}(3,\el)}}\\
& \smashedrightarrow{\delta}
& \cons{1}{\cons{2}{\cons{3}{\el}}}\\
& = & [1;2;3].
\end{array}
\end{equation*}

\end{frame}

% ------------------------------------------------------------------------
%
\begin{frame}{Termination / Insertion sort}

  Informally, what soundness means is that, if some program
  terminates, then the result is what was expected. This property is
  called \textbf{partial correctness} when it is relevant to
  distinguish it from \textbf{total correctness}, which is partial
  correctness and termination.

  \bigskip

  Let us prove the termination of \fun{isort} by the method of
  dependency pairs. The pairs to order are the following:
  \begin{itemize}

    \item \((\fun{ins}(x, \cons{y}{s}), \fun{ins}(x,s))_\gamma\),
    \item \((\fun{isort}(\cons{x}{s}), \fun{isort}\,s)_\beta\) and
    \item \((\fun{isort}(\cons{x}{s}), \fun{ins}(x, \fun{isort}\,s))_\beta\).

  \end{itemize}

\end{frame}

% ------------------------------------------------------------------------
%
\begin{frame}{Termination / Insertion sort}

  By using the proper subterm relation on the second parameter of
  \fun{ins}, we order the first pair:
  \begin{equation*}
    \cons{y}{s} \succ s \Rightarrow \fun{ins}(x, \cons{y}{s}) \succ
    \fun{ins}(x,s).
  \end{equation*}
  This is enough to prove that \fun{ins} terminates. The second pair
  is similarly oriented:
  \begin{equation*}
    \cons{x}{s} \succ s \Rightarrow \fun{isort}(\cons{x}{s}) \succ
    \fun{isort}\,s.
  \end{equation*}

\end{frame}

% ------------------------------------------------------------------------
%
\begin{frame}{Termination / Insertion sort}

  The third pair is not worth considering after all, because we
  already know that \fun{ins} terminates, so the second pair is
  enough to entail the termination of \fun{isort}.

  \bigskip

  In other words, since \fun{ins} terminates, it can be considered,
  as far as termination analysis is concerned, as a data constructor,
  so the third pair becomes:
  \begin{equation*}
    \cons{x}{s} \succ s
    \Rightarrow
    \fun{isort}(\cons{x}{s}) \succ \fun{isort}\,s
    \Rightarrow
    \fun{isort}(\cons{x}{s}) \succ \underline{\fun{ins}}(x, \fun{isort}\,s).
  \end{equation*}
  where \(\fun{\ufun{ins}}\) stands for \fun{ins} considered as a
  \textbf{constructor}, that is, an undefined function.\hfill\(\Box\)

\end{frame}

% ------------------------------------------------------------------------
%
\begin{frame}{Termination / List flattening}

   Let us try a more complicated example:
   \begin{equation*}
     \begin{array}{r@{\;}l@{\;}l}
       \fun{flat}\,\el                   & \smashedrightarrow{\alpha}
       & \el\\
       \fun{flat}(\cons{\el}{t})         & \smashedrightarrow{\beta}
       & \fun{flat}\,t\\
       \fun{flat}(\cons{(\cons{x}{s})}{t}) & \smashedrightarrow{\gamma}
       & \fun{cat}(\fun{flat}(\cons{x}{s}),\fun{flat}\,t)\\
       \fun{flat}(\cons{y}{t})           & \smashedrightarrow{\delta}
       & \cons{y}{\fun{flat}\,t}
     \end{array}
   \end{equation*}

   Some evaluations:
   \begin{align*}
     \fun{flat}\,\el & \twoheadrightarrow \el\\
     \fun{flat}\,[\el;[\el]] & \twoheadrightarrow \el\\
     \fun{flat}\,[\el;[1;[2;\el];3];\el] & \twoheadrightarrow [1;2;3]
   \end{align*}

   What '\fun{flat}' does is flatten a list of lists arbitrarily
   deep.

\end{frame}

% ------------------------------------------------------------------------
%
\begin{frame}{Termination / List flattening}

  A detailed evaluation:
    \begin{equation*}
      \begin{array}{r@{\;}l@{\;}l}
        \fun{flat}\,[\el;[[1];2];3]
        & \xrightarrow{\beta}
        & \fun{flat}\,[[[1];2];3]\\
        & \smashedrightarrow{\gamma}
        & \fun{cat}(\underline{\fun{flat}\,[[1];2]},\fun{flat}\,[3])\\
        & \smashedrightarrow{\gamma}
        & \fun{cat}(\fun{cat}(\underline{\fun{flat}\,[1]},\fun{flat}\,[2]),
        \fun{flat}\,[3])\\
        & \smashedrightarrow{\delta}
        & \fun{cat}(\fun{cat}(\cons{1}{\underline{\fun{flat}\,\el}},
        \fun{flat}\,[2]), \fun{flat}\,[3])\\
        & \smashedrightarrow{\alpha}
        & \fun{cat}(\fun{cat}([1],\underline{\fun{flat}\,[2]}),
        \fun{flat}\,[3])\\
        & \smashedrightarrow{\delta}
        & \fun{cat}(\fun{cat}([1],\cons{2}{\underline{\fun{flat}\,\el}}),
        \fun{flat}\,[3])\\
        & \smashedrightarrow{\alpha}
        & \fun{cat}(\fun{cat}([1],[2]),\underline{\fun{flat}\,[3]})\\
        & \smashedrightarrow{\delta}
        & \fun{cat}(\fun{cat}([1],[2]),\cons{3}{\underline{\fun{flat}\,\el}})\\
        & \smashedrightarrow{\alpha}
        & \fun{cat}(\fun{cat}([1],[2]),[3])\\
        & \twoheadrightarrow & [1;2;3].
      \end{array}
    \end{equation*}

\end{frame}

% ------------------------------------------------------------------------
%
\begin{frame}{Termination / List flattening}

    The function \fun{cat} is independent of \fun{flat} and we
    proved its termination in isolation by using the proper subterm
    order on its first argument.

    \bigskip

    Assuming now that \fun{cat} terminates, let's prove the
    termination of \fun{flat}.

    \bigskip

    Because the recursive calls to \fun{flat} contain only (list)
    constructors, we can try to order their arguments.

    \bigskip

    The proper subterm order works:
    \begin{center}
    \begin{tabular}{r@{\;\,}c@{\;}ll}
      \(\cons{y}{t}\) & \(\succ\) & \(t\) & (rule~\(\delta\)
      and rule~\(\beta\) when \(y=\el\))\\
      \(\cons{(\cons{x}{s})}{t}\) & \(\succ\) & \(t\)
      & (rule~\(\gamma\))\\
      \(\cons{(\cons{x}{s})}{t}\) & \(\succ\) & \(\cons{x}{s}\)
      & (rule~\(\gamma\))
    \end{tabular}
    \end{center}

    \bigskip

    Termination ensues.\hfill\(\Box\)

\end{frame}

% ------------------------------------------------------------------------
%
\begin{frame}{Termination / List flattening}

  Let us consider now an alternative design for rule~\(\gamma\):
  \begin{equation*}
    \begin{array}{@{}r@{\;}c@{\;}l@{}}
      \fun{flat}\,\el & \xrightarrow{\alpha}
      & \el\\
      \fun{flat}(\cons{\el}{t}) & \smashedrightarrow{\beta}
      & \fun{flat}\,t\\
      \fun{flat}(\cons{(\cons{x}{s})}{t})
      & \smashedrightarrow{\gamma}
      & \fbox{\(\fun{flat}(\cons{x}{\cons{s}{t}})\)}\\
      \fun{flat}(\cons{y}{t})   & \smashedrightarrow{\delta}
      & \cons{y}{\fun{flat}\,t}
    \end{array}
  \end{equation*}

  \bigskip

  This amounts to perform a \textbf{right rotation} of the abstract
  syntax tree of the term:
  \begin{center}
    \includegraphics{flat_rotr}
  \end{center}

\end{frame}

% ------------------------------------------------------------------------
%
\begin{frame}{Termination / Measures}

  Here, the proper subterm order fails:
  \begin{equation*}
    \cons{(\cons{x}{s})}{t} \nsucc
    \cons{x}{(\cons{s}{t})} = \cons{x}{\cons{s}{t}}
  \end{equation*}
  because \(t \nsucc \cons{s}{t}\).

  \bigskip

  We need another method to prove termination: \textbf{measures}.

  \bigskip

  A measure \(\measure{\cdot}\) is a map from terms to a
  well\hyp{}ordered set (\(A, \succ\)), which is \textbf{monotonically
    decreasing} with respect to the rewrite relation:
  \begin{equation*}
    x \rightarrow y \Rightarrow \measure{x} \succ_A \measure{y}.
  \end{equation*}

  \medskip

  Actually, as earlier, we will only consider \textbf{dependency
    pairs}, that is, pairs of calls whose first component is the
  left\hyp{}hand side of a rule and the second components are the
  calls in the right\hyp{}hand side of the same rule.

\end{frame}

% ------------------------------------------------------------------------
%
\begin{frame}{Termination / Polynomial measures}

  The pairs are
  \begin{enumerate}

    \item \((\fun{flat}(\cons{\el}{t}), \fun{flat}\,t)_\beta\)

    \item \((\fun{flat}(\cons{(\cons{x}{s})}{t}),
      \fun{flat}(\cons{x}{\cons{s}{t}}))_\gamma\)

    \item \((\fun{flat}(\cons{y}{t}), \fun{flat}\,t)_\delta\)

  \end{enumerate}
  with \(y \not\in S\), that is, \(y\)~is not a list.

  \bigskip

  We can actually drop the function names, as all the pairs involve
  \fun{flat}.

  \bigskip

  A common class of measures are monotonically increasing embeddings
  into (\(\mathbb{N},>\)), so let us seek a measure satisfying:
  \begin{center}
    \begin{tabular}{r@{\;\,}c@{\;}ll}
      \(\measure{\cons{(\cons{x}{s})}{t}}\) & > &
      \(\measure{\cons{x}{\cons{s}{t}}}\)\\
      \(\measure{\cons{y}{t}}\) & > & \(\measure{t}\),
      & if \(y \not\in S\) or \(y=\el\).
    \end{tabular}
  \end{center}

\end{frame}

% ------------------------------------------------------------------------
%
\begin{frame}{Termination / Polynomial measures}

  For instance, let us set the following \textbf{polynomial measure}:
  \begin{center}
    \begin{tabular}{r@{\;\,}c@{\;\;}ll}
      \(\measure{\cons{x}{s}}\) & := & \(1 + 2\cdot\measure{x} +
      \measure{s}\)\\
      \(\measure{y}\) & := & \(0\), & if \(y \not\in S\) or
      \(y=\el\).
    \end{tabular}
  \end{center}

  We have
  \begin{align*}
    \measure{\cons{(\cons{x}{s})}{t}} &  =  3 +
    4\cdot\measure{x} + 2\cdot\measure{s} + \measure{t}\\
    \measure{\cons{x}{\cons{s}{t}}} &  =  2 +
    2\cdot\measure{x} + 2\cdot\measure{s} + \measure{t}
  \end{align*}

  Then
  \begin{center}
    \begin{tabular}{r@{\;\;}c@{\;\;}l}
      \(\measure{\cons{(\cons{x}{s})}{t}}\) & >
      & \(\measure{\cons{x}{\cons{s}{t}}}\)\\
      \(\measure{\cons{y}{t}} = 1 + \measure{t}\) & > &
      \(\measure{t}\)
    \end{tabular}
  \end{center}

  Since \(\measure{x} \in \mathbb{N}\), for all~\(x\), these
  inequalities entail the termination of \fun{flat} (no infinite
  descent).\hfill\(\Box\)

\end{frame}

% ------------------------------------------------------------------------
%
\begin{frame}{Termination / Dependency pairs in more details}

  Let's try to be a little bit more rigorous with dependency pair
  analysis.

  \bigskip

  Perhaps the first thing we have noticed is the use of conditions
  (guards) on rewrite rules. They are actually simple ways to make the
  systems more readable and more efficiently implemented in functional
  languages. In a way, they are akin to macros.

  \bigskip

  For example, consider the definitions of functions calculating the
  maximum and the minimum of two numbers:
  \begin{equation*}
    \begin{array}{@{}r@{\;}l@{\;}lr@{\;}l@{\;}l@{}}
      \fun{max}(0,y) & \rightarrow & y; & \fun{min}(0,y) & \rightarrow & 0;\\
      \fun{max}(x,0) & \rightarrow & x; & \fun{min}(x,0) & \rightarrow & 0;\\
      \fun{max}(x,y) & \rightarrow & 1 + \fun{max}(x\!-\!1,y\!-\!1).
      & \fun{min}(x,y) & \rightarrow & 1 + \fun{min}(x\!-\!1,y\!-\!1).
    \end{array}
  \end{equation*}
  Then insertion sort can be written (very inefficiently) as:
  \begin{equation*}
    \begin{array}{@{}r@{\;}l@{\;}l@{}}
      \fun{ins}(x,\el) & \rightarrow & [x];\\
      \fun{ins}(x,\cons{y}{s}) & \rightarrow &
      \cons{\fun{min}(x,y)}{\fun{ins}(\fun{max}(x,y),s)}.
    \end{array}
  \end{equation*}

\end{frame}

% ------------------------------------------------------------------------
%
\begin{frame}{Termination / Dependency pairs in more details}

  An equivalent but more compact definition would determine maximum
  and the minimum at the same time, like so:
\begin{equation*}
\begin{array}{r@{\;}c@{\;}l}
  \fun{isort}\,\el
& \smashedrightarrow{\alpha}
& \el;\\
  \fun{isort}(\cons{x}{s})
& \smashedrightarrow{\beta}
& \fun{ins}(x,\fun{isort}\,s).\\
  \fun{ins}(x, \el)
& \smashedrightarrow{\gamma}
& [x];\\
   \fun{ins}(x,\cons{y}{s})
& \smashedrightarrow{\delta}
& \fun{choose}(x, \cons{y}{s}, x, y).\\
  \fun{choose}(x, \cons{y}{s}, \fun{Succ}\,a, \fun{Succ}\,b)
& \smashedrightarrow{\epsilon}
& \fun{choose}(x, \cons{y}{s}, a, b);\\
  \fun{choose}(x, \cons{y}{s}, a, 0)
& \smashedrightarrow{\phi}
& \cons{y}{\fun{ins}(x,s)};\\
  \fun{choose}(x, \cons{y}{s}, 0, \fun{Succ}\,b)
& \smashedrightarrow{\psi}
& \cons{x}{\cons{y}{s}}.
\end{array}
\end{equation*}
Here, we use Peano's unary encoding of integers:
\begin{itemize}

  \item the successor function \fun{Succ}, so that \(\fun{Succ}(x)\)
    denotes canonically the value of \(x+1\); (We do not have to assume
    `-'.)

  \item Also, `0' here actually stands for \fun{Zero}().

\end{itemize}

\end{frame}

% ------------------------------------------------------------------------
%
\begin{frame}{Termination / Dependency pairs in more details}

  Instead of denoting the function symbols that are considered
  undefined in a dependency pair, as we did before, we are going to
  mark the defined ones (that is equivalent), with a sharp symbol
  (\#). This enables to write \textbf{the pairs as rewrite rules},
  which is more readable.

  \bigskip

  We have five of them:
  \begin{equation*}
    \begin{array}{r@{\;}c@{\;}ll}
        \fun{isort}^\#(\cons{x}{s})
      & \smashedrightarrow{a}
      & \fun{ins}^\#(x, \fun{isort}\, s), & \text{from \(\beta\)};\\
        \fun{isort}^\#(\cons{x}{s})
      & \smashedrightarrow{b}
      & \fun{isort}^\#\, s, & \text{from \(\beta\)};\\
        \fun{ins}^\#(x,\cons{y}{s})
      & \smashedrightarrow{c}
      & \fun{choose}^\#(x, \cons{y}{s}, x, y), & \text{from \(\delta\)};\\
        \fun{choose}^\#(x, \cons{y}{s}, \fun{Succ}\,a, \fun{Succ}\,b)
      & \smashedrightarrow{d}
      & \fun{choose}^\#(x, \cons{y}{s}, a, b), & \text{from \(\epsilon\)};\\
        \fun{choose}^\#(x, \cons{y}{s}, a, 0)
      & \smashedrightarrow{e}
      & \fun{ins}^\#(x,s), & \text{from \(\phi\)}.
    \end{array}
  \end{equation*}

\end{frame}

% ------------------------------------------------------------------------
%
\begin{frame}{Termination / Dependency pairs in more details}

  We can visualise better the pairs by drawing the dependency graph:
  each vertex is labelled by a pair (\(a\), \(b\), \(c\) etc.), and
  there is an arc \(i \rightarrow j\) if the right\hyp{}hand side of
  pair~\(i\) matches the left\hyp{}hand side of pair~\(j\):
  \begin{center}
    \begin{figure}
      \includegraphics[bb=80 675 268 705]{dp1}
    \end{figure}
  \end{center}
  That graph contains two \textbf{strongly\hyp{}connected components
    (SCCs)}: \(\{b\}\) and \(\{c, d, e\}\). We must show that
  \textbf{all cycles in each SCC map to a strictly decreasing
    measure.}

  \bigskip

  The first one is easy (it is a loop):
  \begin{equation*}
    \cons{x}{s} \succ s
    \Rightarrow
    \fun{isort}^\#(\cons{x}{s}) \smashedrightarrow{b} \fun{isort}^\#\,s
  \end{equation*}
  Another way to put it is that the projection of the pair~\(b\) on
  the parameter of \(\fun{isort}^\#\), that is \(\cons{x}{s}
  \rightarrow s\), if ordered, implies that the original pair is
  ordered too.

\end{frame}

% ------------------------------------------------------------------------
%
\begin{frame}{Termination / Dependency pairs in more details}

  We can annotate the graph with the ordering relation we found:
  \begin{center}
    \begin{figure}
      \includegraphics[bb=80 675 268 705]{dp2}
    \end{figure}
  \end{center}
  For the second and last SCC \(\{c, d, e\}\), if we project
  \(\fun{ins}^\#\) and \(\fun{choose}^\#\) along their second
  parameter, we get:
  \begin{equation*}
    \begin{array}{r@{\;}c@{\;}ll}
      \cons{y}{s} & \smashedrightarrow{c} & \cons{y}{s},\\
      \cons{y}{s} & \smashedrightarrow{d} & \cons{y}{s},\\
      \cons{y}{s} & \smashedrightarrow{e} & s.\\
    \end{array}
  \end{equation*}
  So all paths through~\(e\) strictly decrease the second parameter:
  \begin{center}
    \begin{figure}
      \includegraphics[bb=80 675 268 705]{dp3}
    \end{figure}
  \end{center}

\end{frame}

% ------------------------------------------------------------------------
%
\begin{frame}{Termination / Dependency pairs in more details}

  Now the SCC \(\{c, d, e\}\) reduces to the loop SCC \(\{d\}\).

  \bigskip

  Pair~\(d\) is easy to order: for example, we can project the
  components of the pair along the third parameter of
  \(\fun{choose}^\#\):
  \begin{equation*}
    \fun{Succ} \; a \xrightarrow{d} a.
  \end{equation*}
  So
  \begin{center}
    \begin{figure}
      \includegraphics[bb=80 675 268 705]{dp4}
    \end{figure}
  \end{center}
  (Remember that the immediate subterm ordering \((\succ)\) applies to
  the result of different projections.)

  \bigskip

  This entails that the system is terminating.\hfill$\Box$

\end{frame}
